% October 13, 2026 build (agent draft): conclusion at the 7.9 production
% point. The Oct 12 deck's conclusion frame still says the rebated property
% tax "can move space toward young families"; the 7.9 runs do not support it
% (tax: price -6.2%, births -0.2%, 7p9167_v6/table.md), so it is not repeated.
\section{Conclusion}
\label{sec:conclusion}

I develop a life-cycle model of fertility and housing in
which households choose when to have children, how much space to occupy, and
whether to rent or own it. Children need space, the largest homes must be
bought with a down payment, and earnings are uninsurable. I calibrate the
model to the United States in 2007 and treat the fall in fertility since then
as the start of a transition to a smaller population. As smaller cohorts
enter adulthood, housing demand falls, lowering the cost of space under the
model's supply schedule.

Three results emerge. First, earnings risk matters for fertility: a
reduction in the standard deviation of earnings shocks from
0.48 to 0.35 raises completed fertility by 16 percent. Second, housing costs
matter: a 10 percent expansion of the housing supply raises the completed
fertility of today's young by about 5 percent and of cohorts twenty years
later by about 4 percent. Third, easier mortgage credit mainly changes the
timing of births. It raises births on impact and the ownership of the young,
who buy earlier, smaller and more leveraged homes, but %OLD-20261016 leaves completed
%OLD-20261016 fertility unchanged.
lowers the
completed fertility of today's young slightly (0.4 percent) and raises that of later cohorts by less than 1 percent. %OLD-20261016 Credit tied to children is the exception: restricted to
%OLD-20261016 parents, it raises completed fertility by about 6 percent.
Credit tied to children is the exception: restricted to
parents, it raises completed fertility by about 6 percent, and in the long run
the number of households by about 26 percent, against 2.5 percent for credit
to all. The common thread
is that a child commits a household to paying for family space every period.
Policies raise fertility when they lower that cost or make it easier to bear,
while general mortgage credit primarily changes the route into ownership
in these experiments.

%OLD-20261016 The analysis has several limitations.
Housing policy for births must therefore be judged in the long run and not
only on impact. A policy that raises desired fertility at a given price shows
up in the long run as more households and a higher price, not as more births
per household.

The analysis has several limitations. The moments identify the parameters only
jointly (Section~\ref{sec:quant-identification}): no moment measures how
births respond to the price of space, the response of children ever born by 25
is not matched, and one combination of parameters is nearly unconstrained.
\problem{Preference shocks that affect every living cohort overstate
childlessness and miss the postponement of first births between 2007 and
2023. They also imply a falling house price where the data show a rising one.}
\problem{Lower earnings risk also reduces earnings inequality; the effects
of uncertainty and inequality, and the share operating through housing,
remain to be distinguished.}
The model overstates homeownership among young households, placing too many
of them at the margin where mortgage credit acts. Credit for parents does
not account for default or mortgage-insurance costs. Landlords and builders
are outside the household sector, so income from renting and building housing
does not return to households. These limitations matter for the quantitative
effects of policy and for the interpretation of the fertility decline since
2007.
